\section{本章小结}
\label{sec:rs-summary}

本章讨论只能评价候选时怎样组织搜索。局部改善会停在低谷，这在有限状态的双谷代价和二维多峰函数上都看得很清楚。要越过障碍，算法必须做眼前并不划算的事，四类方法的差别在于保存什么、用什么规则更新。

\textbf{模拟退火}控制单个状态的接受。对称提议下，固定温度的接受规则使Gibbs分布平稳。全局收敛定理只在有限状态、提议连通且$c\geq d^\star$的对数降温下成立，而且是概率的极限。实际的几何降温按有限预算的结果评价，温度要与代价的尺度匹配。

\textbf{演化策略}适应变异分布。保留亲代使最好值不增，整代替换允许当前值上升。CMA-ES用选中后代的加权平均移动中心，用路径和选中位移更新协方差，用标准化路径的长度调整尺度。搜索协方差不是Hessian，也不无条件等于Hessian的逆。高斯搜索分布的均值梯度与第~\ref{chap:zeroth-order-optimization}章的平滑梯度估计重合，排序效用则改变了目标。

\textbf{遗传算法}组织编码的重组与替换。编码决定什么叫邻近，排列上要用保持排列的算子，比例选择要先把代价转换为非负权重。交叉重组已有片段，变异引入新内容，没有变异时种群会完全趋同。

\textbf{粒子群}利用个人和邻域的记忆移动候选。记忆的差不是梯度，没有阻尼时确定性粒子不衰减，速度限幅后要按修正后的速度计算位移。粒子聚在一起不是最优性的证明。

各种方法要落在同一套可行性、评价信息和预算的记账上。不可行候选的处理决定了实际返回的解，修复会改变搜索分布，模拟退火须重新核对接受规则。含噪读数下最低读数偏向有利噪声，对选出的候选要用独立的新读数复核。比较方法时按实际函数调用统一预算，比较多次独立运行的结果分布。

搜索取得改善，不等于找到了全局最优。能给出下界时报告间隙，否则如实报告最好可行值和所用的预算。下一章讨论评价更昂贵的情形，用概率代理模型决定下一次评价的位置。
